\Needspace{18\baselineskip}
\section{Context composition: absence, reuse, and a collision}
\Needspace{24\baselineskip}
\subsection{Hole-free filling inside a box}
The value corolla is hole-free: any hole witness would claim that the lookup
of a primitive operation is a hole label, whereas that lookup is empty.
Applying deep context closure for a hole-free replacement to the old boxed-hole
site proves that its filled result has unique holes. The context-valued operation
also builds a certified filled context from the boxed context and the value
corolla, using hole-free uniqueness and disjointness for the replacement.
\lean{Graph/ContextExamples}{10}{12}{Linear.Examples.value_holeFree}
\lean{Graph/ContextExamples}{14}{15}{Linear.Examples.filledBox_unique}
\lean{Graph/ContextExamples}{17}{19}{Linear.Examples.filledContext}

The single-hole corolla has no retained occurrence when its only outer edge is
removed. A proposed retained witness includes an edge different from zero, but
every element of its edge carrier $\operatorname{Fin}(1)$ equals zero. This
contradicts the witness's inequality certificate.
\lean{Graph/ContextExamples}{21}{25}{Linear.Examples.singleHole_retained_absent}
Consequently replacing its only hole by the very same single-hole corolla still
gives a context. Use sharp context closure: both input uniqueness proofs are
available, and the retained part of the freshness premise is impossible by the
preceding result. However, these two identical input graphs do not have fully
disjoint inventories: their sole edge is a witness for the same label in each.
Applying a proposed full-disjointness proof to those two witnesses gives a
contradiction. Together these examples show that sharp freshness admits valid
label reuse which the stronger sufficient condition would exclude.
\lean{Graph/ContextExamples}{28}{30}{Linear.Examples.reuse_removed_label}
\lean{Graph/ContextExamples}{33}{34}{Linear.Examples.reuse_not_fully_disjoint}

Filling the single-hole corolla with the value is hole-free, not just a context.
The exact inventory theorem leaves a retained alternative or a value alternative;
the first is impossible by absence of surviving old edges and the second by
hole-freeness of the value.
\lean{Graph/ContextExamples}{36}{40}{Linear.Examples.filled_single_holeFree}
The boxed filled result is hole-free at every depth as well. Express any proposed
hole in fibre coordinates. Its outer edge has singleton carrier, so it is zero.
The already proved label equation identifies that edge as a box whose body is
the filled single-hole graph. Transport the fibre address along that equation,
which preserves its hole observation. Its shallow alternative cannot observe a
hole because it is a box; its child alternative contradicts the filled body's
hole-freeness. Thus no shallow or deep witness exists in the result.
\lean{Graph/ContextExamples}{43}{55}{Linear.Examples.filledBox_holeFree}

\Needspace{24\baselineskip}
\subsection{Two distinct shallow holes in series}
For a natural-number name $n$, define a hole label $\chi_n$ with one term input
and one term output. The two-hole graph has external type $*\to *$, one internal
term vertex $v$, and edges $e_0,e_1$ with labels $\chi_0,\chi_1$ and incidence
\[
 \mathsf{input}_0\xrightarrow{e_0}v\xrightarrow{e_1}\mathsf{output}_0.
\]
This notation records the two source and target singleton lists, not an extra
path assumption.
\lean{Graph/ContextExamples}{58}{58}{Linear.Examples.chainHole}
\lean{Graph/ContextExamples}{61}{128}{Linear.Examples.twoHoles}
\begin{proof}[Graph invariants]
The carriers are $\operatorname{Fin}(1)$ and $\operatorname{Fin}(2)$, with finite
range enumerations. Every incidence list is singleton and hence distinct.
For source uniqueness, two edges of different indices have different source
constructors, namely input and internal, so a shared source is impossible;
equal-index cases give equal edges. The target argument uses internal and output
constructors instead. Coverage chooses $e_0$ for the input source and internal
target, and $e_1$ for the internal source and output target. Outer outputs have
no source occurrence and outer inputs no target occurrence by constructor
disjointness. There is only one internal index and one index at each boundary,
so these cases exhaust all vertices. Every displayed vertex label is term,
matching the singleton input and output lists of both typed hole labels.
\end{proof}

This graph has unique holes. Two equal observed labels give equal numeric hole
names. If one edge index is zero and the other is not, their names are respectively
zero and one, a contradiction. If both are zero they agree directly. If neither
is zero, both are one because their ordinal carrier has size two, so they agree
again. At depth zero these are all occurrence addresses.
\lean{Graph/ContextExamples}{130}{149}{Linear.Examples.twoHoles_unique}

\Needspace{24\baselineskip}
\subsection{A checked failure without freshness}
Let the replacement be the corolla of $\chi_1$, with the same external interface
as $e_0$. It is itself a context: any two occurrences belong to its singleton
edge carrier and hence are equal, independently of their lookup hypotheses.
\lean{Graph/ContextExamples}{151}{152}{Linear.Examples.repeatedHole}
\lean{Graph/ContextExamples}{154}{159}{Linear.Examples.repeatedHole_unique}
Nevertheless, replacing $e_0$ in the two-hole graph by this corolla violates
unique holes. In the result, the retained host edge $e_1$ and the inserted
replacement edge both observe $\chi_1$. Context uniqueness would identify their
occurrences, but their edge tags are different constructors of the disjoint sum.
This is a contradiction. Thus well-formed graph plugging and uniqueness of each
input graph alone do not guarantee the context condition in the result.
\lean{Graph/ContextExamples}{162}{167}{Linear.Examples.collision_breaks_unique}
The same witnesses explicitly refute sharp freshness: $e_1$ and its shallow fibre
address witness retained presence of $\chi_1$, and the replacement's only edge
witnesses inserted presence. Feeding both to a proposed freshness proof gives
false. This regression exercises the actual collision excluded by the
necessary-and-sufficient theorem, alongside the valid label-reuse example above.
\lean{Graph/ContextExamples}{169}{172}{Linear.Examples.collision_not_fresh}
